\documentclass[12pt, a4paper]{article}

% --- Packages ---
\usepackage[a4paper, top=2.5cm, bottom=2.5cm, left=3cm, right=2.5cm]{geometry}
\usepackage{graphicx}
\usepackage{amsmath}
\usepackage{amssymb}
\usepackage{booktabs}
\usepackage{array}
\usepackage{longtable}
\usepackage{caption}
\usepackage{float}
\usepackage[hidelinks]{hyperref}
\usepackage{xcolor}
\usepackage{listings}
\usepackage{fancyhdr}
\usepackage{titlesec}
\usepackage{setspace}
\usepackage{parskip}
\usepackage{microtype}
\usepackage{enumitem}
\usepackage[T1]{fontenc}
\usepackage{lmodern}

% --- Page style ---
\pagestyle{fancy}
\fancyhf{}
\fancyhead[L]{\small\textit{Demand Forecasting System}}
\fancyhead[R]{\small\textit{\nouppercase{\leftmark}}}
\fancyfoot[C]{\thepage}
\renewcommand{\headrulewidth}{0.4pt}

% --- Section formatting ---
\titleformat{\section}{\large\bfseries}{\thesection}{1em}{}[\titlerule]
\titleformat{\subsection}{\normalsize\bfseries}{\thesubsection}{1em}{}
\titleformat{\subsubsection}{\normalsize\itshape}{\thesubsubsection}{1em}{}

% --- Code listing style ---
\definecolor{codegray}{rgb}{0.95,0.95,0.95}
\definecolor{keyblue}{rgb}{0.0,0.3,0.7}
\definecolor{commentgreen}{rgb}{0.1,0.5,0.1}
\definecolor{stringred}{rgb}{0.6,0.1,0.1}
\lstset{
    language=Python,
    backgroundcolor=\color{codegray},
    basicstyle=\small\ttfamily,
    keywordstyle=\color{keyblue}\bfseries,
    commentstyle=\color{commentgreen}\itshape,
    stringstyle=\color{stringred},
    showstringspaces=false,
    breaklines=true,
    frame=single,
    framesep=4pt,
    numbers=left,
    numberstyle=\tiny\color{gray},
    numbersep=5pt,
    tabsize=4,
    captionpos=b
}

\setstretch{1.2}
\setlength{\parskip}{6pt}

% =========================================================
\begin{document}

% ===================== COVER PAGE ========================
\begin{titlepage}
    \centering
    \vspace*{1.5cm}

    {\Huge\bfseries Demand Forecasting System\par}
    \vspace{0.4cm}
    {\large\itshape A Multi-Model Ensemble Approach to Retail Demand Prediction\par}

    \vspace{1.5cm}
    \rule{\textwidth}{1pt}
    \vspace{1cm}

    {\large\bfseries Final Project Report\par}

    \vspace{2cm}

    \begin{tabular}{@{}ll@{}}
        \textbf{Course / Programme:} & Bachelor of Computer Science (BCSC) \\[0.3cm]
        \textbf{GitHub Repository:}  & \url{https://github.com/nhbi05/demand-forecasting} \\[0.3cm]
        \textbf{Submission Date:}    & \today \\
    \end{tabular}

    \vspace{0.8cm}
    {\normalsize\bfseries Group Members\par}
    \vspace{0.3cm}

    \renewcommand{\arraystretch}{1.3}
    \begin{tabular}{@{}clc@{}}
        \toprule
        \textbf{No.} & \textbf{Name} & \textbf{Student ID} \\
        \midrule
        1 & Nansereko Housnah       & 24/U/09631/EVE \\
        2 & Suuna Raymond           & 24/U/11403/EVE \\
        3 & Kewoda Joanitah         & 24/U/05727/EVE \\
        4 & Tamale Denis Valelian   & 24/U/11433/EVE \\
        5 & Nabukeera Sumayah Kaswa & 24/U/07899/EVE \\
        \bottomrule
    \end{tabular}

    \vfill
    \rule{\textwidth}{0.5pt}
\end{titlepage}

% ===================== TOC / LISTS =======================
\newpage
\tableofcontents

\newpage
\listoffigures

\listoftables

% ===================== 1. INTRODUCTION ===================
\newpage
\section{Introduction}
\label{sec:introduction}

\subsection{Background}
\label{subsec:background}

Demand forecasting is the process of estimating the quantity of a product that will be sold
over a future time horizon. It underpins a wide range of supply-chain decisions, including
inventory replenishment, staffing, production scheduling, and logistics planning. In a retail
context, accurate demand forecasts reduce stockouts — where insufficient inventory causes
lost sales — and overstock events, where excess inventory ties up capital and drives markdown
losses~\cite{hyndman2021}.

Classical time-series methods such as ARIMA and exponential smoothing have long been applied
to demand forecasting; however, they assume linear relationships and require individual
model fitting per product, limiting scalability and their ability to capture the complex
non-linear seasonal dynamics present in real-world retail data. Machine-learning approaches
have emerged as a compelling alternative: gradient-boosted trees and Random Forests capture
non-linear feature interactions efficiently~\cite{breiman2001}, while deep recurrent
architectures such as LSTM~\cite{hochreiter1997} and GRU~\cite{cho2014} learn sequential
temporal patterns directly from historical observations.

Rather than selecting a single model, ensemble methods combine predictions from multiple
diverse models to reduce variance, smooth individual model weaknesses, and typically
outperform any single constituent model~\cite{dietterich2000}. This project adopts that
ensemble philosophy.

\subsection{Problem Statement}
\label{subsec:problem}

Retail businesses face the challenge of forecasting daily demand at the product level, often
with significant noise in transaction records. No single modelling approach adequately
handles all the variability in real-world sales data: recurrent neural networks capture
sequential dependencies well but may overfit on sparsely observed products, while tree-based
models handle feature interactions effectively but miss long-range temporal patterns.
Probabilistic models such as Gaussian Processes quantify forecast uncertainty but are
computationally prohibitive at scale.

This project addresses these challenges by:
\begin{enumerate}[leftmargin=*, label=(\roman*)]
    \item Designing and implementing a shared preprocessing pipeline that produces canonical
          data artifacts used by all models.
    \item Training four independent forecasting models (LSTM, GRU, Random Forest, Gaussian
          Process) on shared data.
    \item Combining their predictions through a performance-weighted ensemble to produce a
          final forecast.
    \item Evaluating all models and the ensemble on a held-out test set using standard
          regression metrics.
\end{enumerate}

\subsection{Objectives}
\label{subsec:objectives}

The specific objectives of this project are:
\begin{itemize}[leftmargin=*]
    \item Develop a preprocessing pipeline that transforms 7.4~million raw retail
          transactions into a clean, feature-rich daily dataset.
    \item Train LSTM, GRU, Random Forest, and Gaussian Process models on the processed data.
    \item Implement a simple-average ensemble and a performance-weighted ensemble.
    \item Evaluate all models using $R^{2}$, RMSE, and MAPE on a common held-out test set.
    \item Analyse and interpret results to draw conclusions about the relative strengths of
          each modelling approach.
\end{itemize}

\subsection{Significance of the Study}
\label{subsec:significance}

Improved demand forecasting has direct business value: it reduces overstock and stockout
events, informs replenishment decisions, and supports data-driven planning. Beyond the
immediate application, this project demonstrates a replicable pipeline in which multiple
heterogeneous machine learning models are trained independently on shared artifacts and
combined transparently — a pattern applicable across a wide range of time-series
forecasting problems.

% ===================== 2. METHODOLOGY ===================
\newpage
\section{Methodology}
\label{sec:methodology}

% -------- 2.1 Dataset Description --------
\subsection{Dataset Description}
\label{subsec:dataset}

The raw dataset (\texttt{sales.csv}) comprised \textbf{7,432,685} retail transaction rows
spanning \textbf{28~August~2022 to 26~September~2024} (761~calendar days) across
\textbf{4~stores} and \textbf{28,182~unique products}. Each row recorded one transaction
with the six fields described in Table~\ref{tab:raw_schema}.

\begin{table}[H]
    \centering
    \caption{Raw dataset column schema (\texttt{sales.csv}).}
    \label{tab:raw_schema}
    \renewcommand{\arraystretch}{1.2}
    \begin{tabular}{@{}llp{7cm}@{}}
        \toprule
        \textbf{Column} & \textbf{Type} & \textbf{Description} \\
        \midrule
        \texttt{date}        & date   & Transaction date \\
        \texttt{item\_id}    & string & Unique product identifier \\
        \texttt{quantity}    & float  & Units sold; negative values indicate returns \\
        \texttt{price\_base} & float  & Base unit price \\
        \texttt{sum\_total}  & float  & Total revenue
                                        (\texttt{quantity} $\times$ \texttt{price\_base}) \\
        \texttt{store\_id}   & int    & Store identifier (4 stores) \\
        \bottomrule
    \end{tabular}
\end{table}

Table~\ref{fig:raw_sample} shows the first five rows of the raw dataset as loaded from
the notebook.

\begin{figure}[H]
    \centering
    \includegraphics[width=0.92\textwidth]{images/raw_sample.png}
    \caption{First five rows of the raw dataset
             (shape: 7{,}432{,}685 $\times$ 6).}
    \label{fig:raw_sample}
\end{figure}

After aggregation, the processed dataset contained \textbf{76,100 rows} — one row per
product per day — focused on the top~100 products by total positive demand.

% -------- 2.2 Exploratory Data Analysis --------
\subsection{Exploratory Data Analysis}
\label{subsec:eda}

Exploratory analysis was conducted on the raw data to understand demand patterns,
distribution characteristics, and temporal structure before any modelling decisions
were made.

\subsubsection{Quantity Distribution}

Transaction quantities were heavily \textbf{right-skewed}: the median was 2~units, the
mean 5.6~units, and the maximum 4,952~units. Most transactions were small, but a long
tail of bulk orders and high-volume products inflated the mean substantially. This
skewness directly motivated per-product outlier capping and independent product-level
scaling during preprocessing.

\begin{figure}[H]
    \centering
    \includegraphics[width=0.92\textwidth]{images/qty_dist.png}
    \caption{Distribution of transaction quantities (clipped to $[-10,\,100]$ for
             readability), showing the right-skewed demand distribution.}
    \label{fig:qty_dist}
\end{figure}

\subsubsection{Seasonality and Trends}

Figure~\ref{fig:daily_volume} shows aggregate daily sales over the full period, revealing
a gradual upward trend and visible weekly periodicity. Figure~\ref{fig:dow} confirms
strong day-of-week effects. Year-over-year analysis further showed a December demand
spike and a step-up in baseline sales from 2023 to 2024.

\begin{figure}[H]
    \centering
    \includegraphics[width=0.92\textwidth]{images/fig1_daily_volume.png}
    \caption{Daily aggregate transaction volume over the full 761-day dataset period,
             showing a gradual upward trend and weekly periodicity.}
    \label{fig:daily_volume}
\end{figure}

\begin{figure}[H]
    \centering
    \includegraphics[width=0.85\textwidth]{images/fig2_dow_seasonality.png}
    \caption{Average demand by day of week, confirming strong weekly seasonality.}
    \label{fig:dow}
\end{figure}

\subsubsection{Stationarity and Temporal Dependency}

An Augmented Dickey-Fuller (ADF) test returned a statistic of $-0.657$ and a p-value
of $0.858$, confirming the series is \textbf{non-stationary} — consistent with the
visible upward trend. ACF and PACF analysis showed persistent autocorrelation across
all 60~lags and dominant direct effects at lags~1 and~7, establishing weekly demand
history as the key temporal signal.

\begin{figure}[H]
    \centering
    \includegraphics[width=0.92\textwidth]{images/acf_pacf.png}
    \caption{ACF (left) and PACF (right) of aggregate daily quantity over 60 lags.
             Persistent ACF confirms non-stationarity; PACF spikes at lags 1 and 7
             identify the dominant predictors.}
    \label{fig:acf_pacf}
\end{figure}

\subsubsection{Key Findings}

\begin{itemize}[leftmargin=*]
    \item \textbf{Right-skewed demand} --- outlier capping and per-product scaling
          were essential to prevent dominant products from distorting model training.
    \item \textbf{Strong weekly seasonality} --- day-of-week cyclical features were
          included in all models; the GRU uses last week's demand as a residual baseline.
    \item \textbf{Non-stationary trend} --- addressed by the GRU's residual learning
          and the LSTM's learned product embeddings.
    \item \textbf{Lag~1 and lag~7 are the strongest predictors} --- informed the
          60-day lookback window for LSTM/GRU and lag feature selection for Random Forest.
    \item \textbf{High product sparsity} --- the median product had only 72~active days,
          restricting modelling to the top~100 high-volume products.
\end{itemize}

% -------- 2.3 Data Cleaning --------
\subsection{Data Cleaning}
\label{subsec:cleaning}

The preprocessing pipeline (\texttt{src/preprocessing.py}) first cleaned the raw data
before any transformation.

\paragraph{Missing value handling.}
Rows missing a date, item identifier, or quantity were dropped. Missing base prices were
imputed using the per-product median, falling back to the global median and then zero.
Missing store identifiers were assigned a sentinel value
(\texttt{\_\_missing\_store\_\_}). Missing revenue totals were recalculated as
\texttt{quantity} $\times$ \texttt{price\_base}.

\paragraph{Daily aggregation.}
Transaction-level data was aggregated to daily product-level demand. Records were
grouped by \texttt{(date, item\_id)} to produce daily net quantity,
\texttt{positive\_quantity} (sales only), returns volume, transaction count, store count,
mean price, and the binary \texttt{had\_sales} flag. A complete product-date grid was
then constructed so that every product has exactly one row per calendar day.

\paragraph{Outlier handling.}
Per-product outlier caps were computed at the \textbf{99.5th percentile} of the training
distribution for \texttt{quantity}, \texttt{positive\_quantity}, and
\texttt{returns\_quantity}. Values exceeding these caps were winsorised (clipped) to the
cap value. Thresholds were fitted on the training period only to prevent data leakage.

% -------- 2.4 Feature Engineering --------
\subsection{Feature Engineering}
\label{subsec:features}

Three categories of features were constructed from the cleaned daily data.

\paragraph{Cyclical calendar features.}
Sine and cosine transformations were applied to preserve the periodic structure of
calendar variables:
\begin{equation}
    \text{dow\_sin} = \sin\!\left(\frac{2\pi \cdot d_{\text{week}}}{7}\right), \qquad
    \text{dow\_cos} = \cos\!\left(\frac{2\pi \cdot d_{\text{week}}}{7}\right)
    \label{eq:calendar}
\end{equation}
with analogous encodings for day-of-month (period~31) and month (period~12), yielding
6~calendar features in total.

\paragraph{Lag features.}
Lagged demand values at 1, 7, 14, and 30~days prior (\texttt{lag\_1},
\texttt{lag\_7}, \texttt{lag\_14}, \texttt{lag\_30}) were constructed to capture
short-term and medium-term demand history.

\paragraph{Rolling statistics.}
Rolling means and standard deviations over 7-, 14-, and 30-day windows were computed
per product, providing a smoothed view of recent demand level and variability.

% -------- 2.5 Scaling and Normalisation --------
\subsection{Scaling and Normalisation}
\label{subsec:scaling}

A \texttt{MinMaxScaler} was fitted \textbf{per product} on the training period and
applied to seven columns: \texttt{quantity}, \texttt{positive\_quantity},
\texttt{returns\_quantity}, \texttt{price\_base}, \texttt{sum\_total},
\texttt{transaction\_count}, and \texttt{store\_count}. Per-product scaling was chosen
because demand magnitudes vary widely across products; a single global scaler would
compress low-volume products into a near-zero range. Fitting scalers on the training
period only ensured no future information leaked into the transformation.

% -------- 2.6 Data Splitting --------
\subsection{Data Splitting}
\label{subsec:splits}

A \textbf{chronological split} was applied to reflect real forecasting conditions —
the model is always trained on past data and evaluated on future data. All scalers,
outlier thresholds, and product selections were derived exclusively from the training
period. The split boundaries are given in Table~\ref{tab:splits}.

\begin{table}[H]
    \centering
    \caption{Chronological train / validation / test split boundaries.}
    \label{tab:splits}
    \renewcommand{\arraystretch}{1.2}
    \begin{tabular}{@{}lcc@{}}
        \toprule
        \textbf{Split} & \textbf{Start Date} & \textbf{End Date} \\
        \midrule
        Training   & 28 Aug 2022 & 27 Feb 2024 \\
        Validation & 28 Feb 2024 & 27 Apr 2024 \\
        Test       & 28 Apr 2024 & 26 Sep 2024 \\
        \bottomrule
    \end{tabular}
\end{table}

Using a time-based split — rather than a random split — prevents data leakage and
ensures that evaluation metrics reflect the model's ability to generalise to genuinely
unseen future demand.

\subsection{System Architecture}
\label{subsec:architecture}

The system was structured in two layers: a shared preprocessing pipeline and four
independent model notebooks, all feeding into a single ensemble notebook. The preprocessing
module produced a set of canonical shared artifacts — \texttt{daily.csv},
\texttt{splits.json}, \texttt{selected\_products.json}, \texttt{scalers.pkl},
\texttt{outlier\_caps.json}, \texttt{feature\_columns.json}, and \texttt{metadata.json}
— that every model notebook consumed without modification. This design guaranteed that
product selection, date splits, scaling, and feature definitions were identical across all
models, making their predictions directly comparable and ensemble-combinable.

Each model notebook saved test-set predictions as a CSV with the schema
\texttt{(date, item\_id, predicted\_quantity, actual\_quantity)}. The ensemble notebook
(\texttt{06\_Ensemble.ipynb}) joined these four files on the composite key
\texttt{(date, item\_id)} and computed weighted averages.

\subsection{Model Descriptions}
\label{subsec:models}

\subsubsection{Long Short-Term Memory (LSTM)}
\label{ssubsec:lstm}

LSTM is a recurrent neural network designed for sequential learning~\cite{hochreiter1997}.
It uses a gating mechanism comprising input, forget, and output gates to selectively retain
or discard information across long sequences:
\begin{align}
    f_t &= \sigma(W_f [h_{t-1}, x_t] + b_f), \label{eq:forget}\\
    i_t &= \sigma(W_i [h_{t-1}, x_t] + b_i), \label{eq:input}\\
    \tilde{c}_t &= \tanh(W_c [h_{t-1}, x_t] + b_c), \label{eq:cell_cand}\\
    c_t &= f_t \odot c_{t-1} + i_t \odot \tilde{c}_t, \label{eq:cell}\\
    o_t &= \sigma(W_o [h_{t-1}, x_t] + b_o), \label{eq:output}\\
    h_t &= o_t \odot \tanh(c_t). \label{eq:hidden}
\end{align}

The model was implemented in PyTorch and accepts a \textbf{60-day lookback window} of
15~features (6~cyclical calendar, 6~scaled demand/price signals, and 3~engineered lag and
rolling features). It additionally learns a \textbf{product embedding} of dimension~4 for
each of the 100~products, implicitly capturing product-level demand characteristics. The
output is a 30-day-ahead forecast vector. Training used Huber loss (SmoothL1Loss), the
Adam optimiser ($\eta = 10^{-3}$, weight decay $= 10^{-4}$), and early stopping with
patience~5 on the validation loss (best epoch:~3, best validation Huber
loss:~0.0135).

\subsubsection{Gated Recurrent Unit (GRU)}
\label{ssubsec:gru}

GRU is a simplified recurrent architecture with two gates (reset and update) instead of
three~\cite{cho2014}, yielding fewer parameters and faster convergence:
\begin{align}
    r_t &= \sigma(W_r [h_{t-1}, x_t]), \label{eq:reset}\\
    z_t &= \sigma(W_z [h_{t-1}, x_t]), \label{eq:update}\\
    \tilde{h}_t &= \tanh(W[r_t \odot h_{t-1},\, x_t]), \label{eq:cand_h}\\
    h_t &= (1 - z_t) \odot h_{t-1} + z_t \odot \tilde{h}_t. \label{eq:gru_h}
\end{align}

A key design choice is \textbf{residual learning}: the GRU predicts a correction
$\Delta_t$ on top of a seasonal-naive baseline (last week's demand), so that
\begin{equation}
    \hat{y}_t = y_{t-7} + \Delta_t.
    \label{eq:residual}
\end{equation}
This focuses model capacity on the portion of demand that weekly seasonality alone cannot
explain. The GRU shares the same architectural hyperparameters as the LSTM
(hidden size~64, 2~layers, dropout~0.3) and uses 14~input features. Training used MSE loss
on the residual target (best epoch:~9).

\subsubsection{Random Forest}
\label{ssubsec:rf}

Random Forest is an ensemble of decision trees trained on bootstrap samples of the
data~\cite{breiman2001}. Unlike the recurrent models it operates on \textbf{flat feature
vectors}. For each (product, day) pair, 21~features are constructed:
\begin{itemize}[leftmargin=*, nosep]
    \item Lag features: demand from 1, 7, 14, and 30~days prior.
    \item Rolling statistics: 7-, 14-, and 30-day rolling means and standard deviations
          (6~features).
    \item Calendar features: 6~cyclical encodings of day-of-week, day-of-month, and month.
    \item Demand and price signals: scaled \texttt{quantity}, \texttt{positive\_quantity},
          \texttt{price\_base}, \texttt{store\_count}, \texttt{had\_sales}.
\end{itemize}
The model used 100~estimators, maximum depth~15, minimum 5~samples to split a node, and
minimum 2~samples per leaf. Predictions for the 30-day horizon were generated
\emph{recursively}, with prior predictions substituted as lag inputs where required.
Feature importance was assessed using the built-in Gini Impurity (Mean Decrease in
Impurity).

\subsubsection{Gaussian Process}
\label{ssubsec:gp}

Gaussian Process (GP) regression is a non-parametric probabilistic model that places a
prior distribution over functions and updates it with observed data~\cite{rasmussen2006}:
\begin{equation}
    f(\mathbf{x}) \sim \mathcal{GP}\!\left(m(\mathbf{x}),\; k(\mathbf{x}, \mathbf{x}')\right).
    \label{eq:gp}
\end{equation}
The GP produces both a \textbf{mean prediction} and a \textbf{predictive standard
deviation}, enabling probabilistic inventory decisions such as safety-stock sizing. The
composite kernel used was:
\begin{equation}
    k = k_{\mathrm{RBF}} + k_{\mathrm{ExpSine}} + k_{\mathrm{Const}} + k_{\mathrm{White}},
    \label{eq:kernel}
\end{equation}
capturing smooth trends (RBF), seasonal patterns (ExpSineSquared), a global constant
shift (ConstantKernel), and observation noise (WhiteKernel).

GP inference scales as $O(n^{3})$, requiring a hard cap of \textbf{200 rows per product}
to keep training tractable. The 11~input features comprised day-of-year and day-of-week
sine/cosine encodings, month encoding, 1-day and 7-day lag features, and 7- and 30-day
rolling means.

\subsection{Ensemble Strategy}
\label{subsec:ensemble}

The ensemble combined the four models' test-set predictions through a weighted average.
Let $\hat{y}^{(m)}_{t}$ denote model~$m$'s prediction at time~$t$ and $w_m$ its weight.
The ensemble forecast is:
\begin{equation}
    \hat{y}^{\mathrm{ens}}_{t} = \sum_{m \in \mathcal{M}} w_m\, \hat{y}^{(m)}_{t},
    \qquad \sum_{m} w_m = 1.
    \label{eq:ensemble}
\end{equation}

Two strategies were implemented:
\begin{enumerate}[leftmargin=*, label=(\roman*)]
    \item \textbf{Simple Average:} $w_m = 0.25$ for all models.
    \item \textbf{Weighted Average:} weights proportional to each model's calibration-half
          $R^{2}$:
          \begin{equation}
              w_m = \frac{\max(R^{2}_m,\, 0)}{\displaystyle\sum_{m'} \max(R^{2}_{m'},\, 0)}.
              \label{eq:weights}
          \end{equation}
\end{enumerate}

To avoid data leakage, weights were calibrated on the \textbf{first half} of the test
window (28~April -- 11~July~2024) and ensemble metrics reported exclusively on the
\textbf{second half} (12~July -- 24~September~2024). The calibrated weights are given in
Table~\ref{tab:weights}.

\begin{table}[H]
    \centering
    \caption{Calibrated ensemble weights (proportional to calibration-half $R^{2}$).}
    \label{tab:weights}
    \renewcommand{\arraystretch}{1.2}
    \begin{tabular}{@{}lc@{}}
        \toprule
        \textbf{Model} & \textbf{Weight} \\
        \midrule
        LSTM             & 0.2551 \\
        GRU              & 0.2549 \\
        Random Forest    & 0.2530 \\
        Gaussian Process & 0.2371 \\
        \bottomrule
    \end{tabular}
\end{table}

The near-equal weights reflect the similar $R^{2}$ performance of all four models on the
calibration window. The GP received a marginally lower weight, consistent with its weaker
point-forecast performance.

\subsection{Technologies Used}
\label{subsec:technologies}

The system was implemented in Python~3, using the tools listed in
Table~\ref{tab:tech}.

\begin{table}[H]
    \centering
    \caption{Technologies and libraries used in this project.}
    \label{tab:tech}
    \renewcommand{\arraystretch}{1.2}
    \begin{tabular}{@{}ll@{}}
        \toprule
        \textbf{Category} & \textbf{Tool / Library} \\
        \midrule
        Deep learning     & PyTorch~\cite{paszke2019} \\
        Traditional ML    & scikit-learn~\cite{pedregosa2011} \\
        Data processing   & pandas, NumPy \\
        Visualisation     & matplotlib \\
        Notebooks         & Jupyter \\
        Language          & Python 3 \\
        \bottomrule
    \end{tabular}
\end{table}

% ===================== 3. RESULTS ========================
\newpage
\section{Results and Discussion}
\label{sec:results}

All evaluation metrics were computed on the \textbf{second half of the test window}
(12~July -- 24~September~2024, 5,775~row-days across the 77~products present in all four
prediction files) unless stated otherwise. Predictions are reported in original demand
units (quantities), not scaled values. MAPE was computed only on rows where actual demand
was strictly positive to avoid division by zero.

\subsection{Individual Model Performance}
\label{subsec:individual}

Table~\ref{tab:individual} reports the performance of each individual model on the
evaluation half of the test window.

\begin{table}[H]
    \centering
    \caption{Individual model performance on the evaluation half of the test set
             (12~Jul -- 24~Sep~2024).}
    \label{tab:individual}
    \renewcommand{\arraystretch}{1.3}
    \begin{tabular}{@{}lcccc@{}}
        \toprule
        \textbf{Model} & \textbf{$R^{2}$} & \textbf{RMSE} & \textbf{MAE} & \textbf{MAPE} \\
        \midrule
        LSTM             & 0.9813 & 64.29  & 35.07 & 37.04\% \\
        GRU              & 0.9806 & 65.61  & 34.79 & 38.45\% \\
        Random Forest    & 0.9664 & 86.20  & 28.97 & 16.54\% \\
        Gaussian Process & 0.8656 & 172.54 & 61.13 & 35.77\% \\
        \bottomrule
    \end{tabular}
\end{table}

\begin{figure}[H]
    \centering
    \includegraphics[width=0.92\textwidth]{images/lstm_pred.png}
    \caption{Actual versus predicted daily demand for LSTM on the evaluation half of the
             test window.}
    \label{fig:lstm_pred}
\end{figure}

\begin{figure}[H]
    \centering
    \includegraphics[width=0.92\textwidth]{images/gru_pred.png}
    \caption{Actual versus predicted daily demand for GRU on the evaluation half of the
             test window.}
    \label{fig:gru_pred}
\end{figure}

\begin{figure}[H]
    \centering
    \includegraphics[width=0.92\textwidth]{images/rf_pred.png}
    \caption{Actual versus predicted daily demand for Random Forest on the evaluation half
             of the test window.}
    \label{fig:rf_pred}
\end{figure}

\subsubsection{LSTM and GRU}

LSTM and GRU delivered the strongest overall performance, with $R^{2}$ of 0.9813 and
0.9806 and RMSE of 64.29 and 65.61 respectively. Their near-identical scores reflect the
well-documented convergence of these two architectures on well-structured time series: the
GRU's simpler gating scheme is sufficient to capture the dominant weekly periodicities and
trends present in retail demand. Both models outperformed the seasonal-naive baseline
(RMSE~54.73 on a 1-day-ahead task) while forecasting 30~days ahead, a substantially harder
horizon. Their higher MAPE values (37--38\%) compared to Random Forest occur because they
attempt to forecast all days — including near-zero-demand periods — where absolute
percentage errors amplify.

\subsubsection{Random Forest}

Random Forest achieved $R^{2} = 0.9664$ with a notably lower MAPE of \textbf{16.54\%},
outperforming all other models on that metric. This suggests strong calibration on
typical-demand days, where short-term lag features and rolling averages are highly
predictive. Its higher RMSE (86.20) relative to the deep learning models indicates that it
handles extreme demand spikes less gracefully, since decision trees cannot smoothly
extrapolate beyond the training range.

\subsubsection{Gaussian Process}

Gaussian Process recorded the weakest point-prediction metrics ($R^{2} = 0.8656$, RMSE
$= 172.54$). This is consistent with two structural constraints: (i) as a probabilistic
model, its primary design goal is calibrated uncertainty rather than minimum point-forecast
error; and (ii) training was capped at 200~rows per product to manage $O(n^{3})$
complexity. Despite its weaker accuracy, the GP uniquely provides a predictive confidence
interval, which is directly applicable to safety-stock calculations in inventory
management.

\subsection{Feature Importance}
\label{subsec:importance}

Feature importance was assessed via the Random Forest's built-in Gini Impurity (Mean
Decrease in Impurity).

\begin{figure}[H]
    \centering
    \includegraphics[width=0.92\textwidth]{images/feature_importance.png}
    \caption{Top-10 feature importances from the Random Forest model assessed by Gini
             Impurity.}
    \label{fig:importance}
\end{figure}

Short-term lag features — particularly \texttt{lag\_1} (previous day's demand) and
\texttt{lag\_7} (seven-day lagged demand) — accounted for the largest share of predictive
importance, confirming that recent sales history and weekly periodicity are the dominant
signals in retail demand. Rolling means over 7 and 14~days ranked next, followed by
cyclical calendar encodings (\texttt{dow\_sin}/\texttt{dow\_cos}) and price signals. This
ranking is consistent with the strong weekly seasonality visible in Figure~\ref{fig:dow}
and validates the feature engineering choices made across all four models.

\subsection{Ensemble Performance}
\label{subsec:ensemble_results}

Both ensemble strategies were evaluated on the second half of the test window after
calibrating weights on the first half. Results are shown in Table~\ref{tab:ensemble}.

\begin{table}[H]
    \centering
    \caption{Ensemble versus individual model performance on the evaluation half of the
             test set (12~Jul -- 24~Sep~2024). Best values in bold.}
    \label{tab:ensemble}
    \renewcommand{\arraystretch}{1.3}
    \begin{tabular}{@{}lcccc@{}}
        \toprule
        \textbf{Model} & \textbf{$R^{2}$} & \textbf{RMSE} & \textbf{MAE} & \textbf{MAPE} \\
        \midrule
        LSTM             & \textbf{0.9813} & \textbf{64.29} & 35.07 & 37.04\% \\
        GRU              & 0.9806          & 65.61          & 34.79 & 38.45\% \\
        Random Forest    & 0.9664          & 86.20          & \textbf{28.97} & \textbf{16.54\%} \\
        Gaussian Process & 0.8656          & 172.54         & 61.13 & 35.77\% \\
        \midrule
        Simple Ensemble  & 0.9670          & 85.43          & 34.04 & 28.09\% \\
        Weighted Ensemble & 0.9681         & 84.06          & 33.70 & 28.08\% \\
        \bottomrule
    \end{tabular}
\end{table}

\begin{figure}[H]
    \centering
    \includegraphics[width=0.92\textwidth]{images/ensemble_comparison.png}
    \caption{Comparison of $R^{2}$, RMSE, and MAPE across all four individual models and
             both ensemble strategies.}
    \label{fig:ensemble_comparison}
\end{figure}

The weighted ensemble marginally outperformed the simple average across all three metrics
($R^{2}$: 0.9681 vs 0.9670; RMSE: 84.06 vs 85.43; MAPE: 28.08\% vs 28.09\%), confirming
that performance-proportional weighting provides a consistent, if modest, benefit. Both
ensembles achieved a MAPE of approximately 28\%, a meaningful improvement over the deep
learning models (37--38\%) while falling between the deep learning and Random Forest
extremes.

Neither ensemble surpassed LSTM or GRU individually on $R^{2}$ and RMSE. This is
consistent with ensemble theory: when the two strongest constituent models are already
highly correlated — both are recurrent sequence models with identical hyperparameters —
variance reduction is limited, and the inclusion of the weaker Gaussian Process introduces
a downward pull on overall accuracy. In settings where constituent models are more
architecturally diverse, ensemble gains are typically larger.

\subsection{Key Insights}
\label{subsec:insights}

\begin{enumerate}[leftmargin=*, label=(\arabic*)]
    \item \textbf{Sequential models dominate accuracy.} LSTM and GRU's ability to learn
          temporal dependencies over a 60-day window makes them the strongest
          point-forecasters for daily retail demand.
    \item \textbf{Random Forest provides complementary calibration.} Its lower MAPE makes
          it the best choice for predicting typical-demand days, where lag and rolling
          features are highly informative.
    \item \textbf{Gaussian Process contributes uniquely through uncertainty.} Despite weaker
          point accuracy, its probabilistic output is directly applicable to risk-aware
          inventory planning.
    \item \textbf{Ensemble gains are modest when models are correlated.} The near-equal
          weights and limited $R^{2}$ improvement reflect the high architectural similarity
          between LSTM and GRU. True ensemble benefit requires model diversity.
    \item \textbf{Feature engineering validates domain knowledge.} The dominance of lag
          features in Random Forest importance confirms that recent demand history is the
          single most predictive signal, consistent with prior literature~\cite{hyndman2021}.
\end{enumerate}

% ===================== 4. CONCLUSION =====================
\newpage
\section{Conclusion}
\label{sec:conclusion}

This project successfully developed a multi-model demand forecasting system trained on
\textbf{7,432,685} real-world retail transactions. Four machine learning architectures ---
LSTM, GRU, Random Forest, and Gaussian Process --- were trained independently on shared
preprocessed data and combined through a weighted ensemble, meeting all stated objectives.

Among individual models, LSTM and GRU achieved the highest predictive accuracy
($R^{2} = 0.981$), demonstrating that deep learning architectures are highly effective at
capturing the sequential, cyclical patterns of daily retail demand. Random Forest delivered
the lowest MAPE (16.54\%), making it the most reliable model for typical-demand days.
Gaussian Process contributed probabilistic uncertainty estimates not available from the
other models.

The weighted ensemble ($R^{2} = 0.968$, RMSE\,=\,84.06, MAPE\,=\,28.08\%) improved
prediction stability and reduced MAPE relative to the deep learning models alone,
demonstrating the value of combining heterogeneous models. The system provides a robust and
scalable framework for retail demand forecasting, equipping decision-makers with daily,
product-level forecasts that can reduce stockouts, optimise replenishment schedules, and
improve service levels and profitability.

% ===================== REFERENCES ========================
\newpage
\begin{thebibliography}{9}

\bibitem{breiman2001}
Breiman, L. (2001).
Random Forests.
\textit{Machine Learning}, 45(1), 5--32.

\bibitem{cho2014}
Cho, K., van Merrienb\"{o}er, B., Gulcehre, C., Bahdanau, D., Bougares, F., Schwenk, H.,
and Bengio, Y. (2014).
Learning Phrase Representations using RNN Encoder-Decoder for Statistical Machine
Translation.
In \textit{Proceedings of the 2014 Conference on Empirical Methods in Natural Language
Processing (EMNLP)}, pp.\ 1724--1734.

\bibitem{dietterich2000}
Dietterich, T.~G. (2000).
Ensemble Methods in Machine Learning.
In \textit{Multiple Classifier Systems}, Lecture Notes in Computer Science, vol.\ 1857.
Springer, pp.\ 1--15.

\bibitem{hochreiter1997}
Hochreiter, S. and Schmidhuber, J. (1997).
Long Short-Term Memory.
\textit{Neural Computation}, 9(8), 1735--1780.

\bibitem{hyndman2021}
Hyndman, R.~J. and Athanasopoulos, G. (2021).
\textit{Forecasting: Principles and Practice}, 3rd~ed.
OTexts.

\bibitem{paszke2019}
Paszke, A. et al.\ (2019).
PyTorch: An Imperative Style, High-Performance Deep Learning Library.
In \textit{Advances in Neural Information Processing Systems}, 32.

\bibitem{pedregosa2011}
Pedregosa, F. et al.\ (2011).
Scikit-learn: Machine Learning in Python.
\textit{Journal of Machine Learning Research}, 12, 2825--2830.

\bibitem{rasmussen2006}
Rasmussen, C.~E. and Williams, C.~K.~I. (2006).
\textit{Gaussian Processes for Machine Learning}.
MIT Press.

\end{thebibliography}

% ===================== APPENDIX A: HYPERPARAMS ===========
\newpage
\appendix

\section{Model Hyperparameters}
\label{app:hyperparams}

\subsection{LSTM}

\begin{table}[H]
    \centering
    \caption{LSTM hyperparameters.}
    \label{tab:lstm_hp}
    \renewcommand{\arraystretch}{1.2}
    \begin{tabular}{@{}ll@{}}
        \toprule
        \textbf{Hyperparameter} & \textbf{Value} \\
        \midrule
        Framework                & PyTorch \\
        Lookback window          & 60 days \\
        Forecast horizon         & 30 days \\
        Hidden units             & 64 \\
        Recurrent layers         & 2 \\
        Dropout rate             & 0.3 \\
        Product embedding dim    & 4 \\
        Optimiser                & Adam ($\eta = 10^{-3}$, weight decay $= 10^{-4}$) \\
        Loss function            & Huber (SmoothL1Loss) \\
        Max epochs               & 100 \\
        Early stopping patience  & 5 (monitor: validation Huber loss) \\
        Best epoch               & 3 \\
        Best validation Huber loss & 0.013505 \\
        Input features           & 15 \\
        \bottomrule
    \end{tabular}
\end{table}

\subsection{GRU}

\begin{table}[H]
    \centering
    \caption{GRU hyperparameters.}
    \label{tab:gru_hp}
    \renewcommand{\arraystretch}{1.2}
    \begin{tabular}{@{}ll@{}}
        \toprule
        \textbf{Hyperparameter} & \textbf{Value} \\
        \midrule
        Framework                & PyTorch \\
        Lookback window          & 60 days \\
        Forecast horizon         & 30 days \\
        Hidden units             & 64 \\
        Recurrent layers         & 2 \\
        Dropout rate             & 0.3 \\
        Product embedding dim    & 4 \\
        Optimiser                & Adam ($\eta = 10^{-3}$, weight decay $= 10^{-4}$) \\
        Loss function            & MSE (on residual target $\Delta_t$) \\
        Max epochs               & 100 \\
        Early stopping patience  & 10 (monitor: validation MSE) \\
        Best epoch               & 9 \\
        Residual learning        & Seasonal-naive baseline $y_{t-7}$ \\
        Seasonality period       & 7 days \\
        Input features           & 14 \\
        \bottomrule
    \end{tabular}
\end{table}

\subsection{Random Forest}

\begin{table}[H]
    \centering
    \caption{Random Forest hyperparameters.}
    \label{tab:rf_hp}
    \renewcommand{\arraystretch}{1.2}
    \begin{tabular}{@{}ll@{}}
        \toprule
        \textbf{Hyperparameter} & \textbf{Value} \\
        \midrule
        Framework               & scikit-learn \\
        \texttt{n\_estimators}  & 100 \\
        \texttt{max\_depth}     & 15 \\
        \texttt{min\_samples\_split} & 5 \\
        \texttt{min\_samples\_leaf}  & 2 \\
        Lag features            & 1, 7, 14, 30 days \\
        Rolling windows         & 7, 14, 30 days (mean + std) \\
        Forecast strategy       & Recursive (30-step) \\
        Random seed             & 42 \\
        Input features          & 21 \\
        Training rows           & 51,900 \\
        \bottomrule
    \end{tabular}
\end{table}

\subsection{Gaussian Process}

\begin{table}[H]
    \centering
    \caption{Gaussian Process hyperparameters.}
    \label{tab:gp_hp}
    \renewcommand{\arraystretch}{1.2}
    \begin{tabular}{@{}ll@{}}
        \toprule
        \textbf{Hyperparameter} & \textbf{Value} \\
        \midrule
        Framework                  & scikit-learn \\
        Kernel                     & RBF + ConstantKernel + WhiteKernel + ExpSineSquared \\
        Max training rows/product  & 200 ($O(n^{3})$ constraint) \\
        Optimiser restarts         & 2 \\
        Output                     & Mean prediction + uncertainty (std) \\
        Input features             & 11 \\
        \bottomrule
    \end{tabular}
\end{table}

% ===================== APPENDIX B: CODE ==================
\newpage
\section{Key Code Snippets}
\label{app:code}

\subsection{Ensemble Prediction Join}

Predictions from the four model files are joined on the composite key
\texttt{(date, item\_id)} before averaging, guaranteeing correct product-day alignment.

\begin{lstlisting}[caption={Inner-join of all four model prediction files.},
                   label=lst:join]
merged = (
    lstm_df
    .rename(columns={'predicted_quantity': 'lstm_pred',
                     'actual_quantity':    'actual'})
    .merge(
        gru_df.rename(columns={'predicted_quantity': 'gru_pred'})
              [['date', 'item_id', 'gru_pred']],
        on=['date', 'item_id'])
    .merge(
        rf_df.rename(columns={'predicted_quantity': 'rf_pred'})
             [['date', 'item_id', 'rf_pred']],
        on=['date', 'item_id'])
    .merge(
        gaussian_df.rename(columns={'predicted_quantity': 'gaussian_pred'})
                   [['date', 'item_id', 'gaussian_pred']],
        on=['date', 'item_id'])
    .sort_values(['date', 'item_id'])
    .reset_index(drop=True)
)
\end{lstlisting}

\subsection{MAPE Calculation}

MAPE is computed only on rows with strictly positive actual demand to prevent division by
zero.

\begin{lstlisting}[caption={MAPE with zero-demand masking.}, label=lst:mape]
def mape(y_true, y_pred):
    mask   = y_true > 0
    y_true = y_true[mask]
    y_pred = y_pred[mask]
    return np.mean(np.abs((y_true - y_pred) / y_true)) * 100
\end{lstlisting}

\subsection{Weighted Ensemble Calibration}

Weights are calibrated on the first half of the test window; the ensemble is then
evaluated exclusively on the second half.

\begin{lstlisting}[caption={Leakage-free weighted ensemble calibration and evaluation.},
                   label=lst:weights]
mid   = len(merged) // 2
calib = merged.iloc[:mid]   # 28 Apr -- 11 Jul 2024
eval_ = merged.iloc[mid:]   # 12 Jul -- 24 Sep 2024

weights       = compute_r2_weights(calib)
ensemble_pred = weighted_average(eval_, weights)
results       = evaluate(eval_['actual_quantity'], ensemble_pred)
\end{lstlisting}

\subsection{Cyclical Calendar Feature Engineering}

Sine/cosine encodings preserve the cyclic structure of calendar variables.

\begin{lstlisting}[caption={Cyclical calendar feature generation
                   (\texttt{src/preprocessing.py}).}, label=lst:calendar]
def add_calendar_features(df):
    df['dow_sin']   = np.sin(2 * np.pi * df['dayofweek'] / 7)
    df['dow_cos']   = np.cos(2 * np.pi * df['dayofweek'] / 7)
    df['dom_sin']   = np.sin(2 * np.pi * df['day']       / 31)
    df['dom_cos']   = np.cos(2 * np.pi * df['day']       / 31)
    df['month_sin'] = np.sin(2 * np.pi * df['month']     / 12)
    df['month_cos'] = np.cos(2 * np.pi * df['month']     / 12)
    return df
\end{lstlisting}

\end{document}
